% DataDiffFuzz — ICSE/FSE submission skeleton.
% Build: pdflatex main && bibtex main && pdflatex main && pdflatex main
% Requires the ACM acmart class (TeX Live `acmart`, or Overleaf).
% NOTE: no LaTeX toolchain is installed in the working environment; verify on Overleaf/TL.
\documentclass[sigconf,review,anonymous]{acmart}

\usepackage{booktabs}
\usepackage{graphicx}
\usepackage{amsmath}
\usepackage{xspace}
\usepackage{enumitem}

\newcommand{\todo}[1]{\textcolor{red}{[TODO: #1]}}
\newcommand{\sys}{\textsc{DataDiffFuzz}\xspace}
\newcommand{\osc}{\textsc{Osc}\xspace}

\settopmatter{printacmref=false}
\renewcommand\footnotetextcopyrightpermission[1]{}
\pagestyle{fancy}
\fancyhf{}

\begin{document}

\title{\sys{}: Unified Semantic Differential Fuzzing for Heterogeneous
Tabular Data-Processing Systems}

\author{Anonymous}
\affiliation{\institution{Under review}\country{}}

\begin{abstract}
Modern tabular data-processing stacks combine Python DataFrame APIs, Arrow-based in-memory
processing, embedded SQL engines, and analytical query engines. These systems are routinely
treated as interchangeable executors of the same table semantics, yet they silently diverge on
null and negative-zero handling, ordering and top-$k$ ties, group-by and join semantics, type
casts, Arrow memory layout, and optimizer rewrites. Existing wrong-result testing techniques
target SQL DBMSs or a single data model, leaving this heterogeneous ecosystem without a common
testing surface. We present \sys{}, a unified semantic differential fuzzing framework for
tabular data-processing systems. \sys{} uses a typed workflow IR with capability-aware lowering
to multiple backend families, a shared semantic normalizer for cross-backend canonicalization,
and a compositional oracle complex that fuses differential checking, metamorphic relations, and
deterministic latest-version probes under a single finite multi-endpoint contract. An end-to-end
evidence pipeline carries each finding through candidate triage, recheck, reduction,
family-level deduplication, issue bundling, and confirmation tracking, keeping fresh discovery
physically separate from replayed and already-known roots. Across DataFusion, Polars, DuckDB,
and PyArrow we report nine upstream-confirmed latest-version bug families (plus one pending and
one non-original match shown separately), and we evaluate noise control, oracle complementarity,
scheduling efficiency, evidence actionability, and cross-family transferability. \todo{replace
with the regenerated result sentence after RQ1--RQ6}.
\end{abstract}

\keywords{differential testing, metamorphic testing, fuzzing, DataFrame systems, query engines,
silent wrong-result bugs}

\maketitle

\section{Introduction}
\label{sec:intro}

\subsection{Setting}

Data-processing pipelines rarely stay inside one engine. A single analyst workflow may read
Parquet through PyArrow, reshape it with pandas or Polars, push an aggregation into DuckDB or
SQLite, and execute a larger relational plan in DataFusion or chDB. These systems are commonly
treated as interchangeable executors of the same table semantics, and in the common case they
behave as such. At the boundaries, however, they disagree: null and NaN propagation,
negative-zero comparisons, ordering and top-$k$ tie-breaking, group-by with null keys, join
semantics under duplicate keys, integer/float casting, string collation, and Arrow layout
details such as slice offsets and chunking. When such a disagreement occurs, both engines
return a plausible result and no error is raised. The pipeline continues with silently wrong
data.

\subsection{Problem}

Detecting these \emph{silent wrong-result} bugs in this ecosystem is hard for four reasons.
First, the testing surface is heterogeneous: one workflow must be lowered into DataFrame API
calls, Arrow compute kernels, SQL, and query-engine plans, each with its own type system and
capability set. Second, directly comparing raw outputs is dominated by false positives, because
legitimate differences in dtype, index, ordering, and presentation are ubiquitous; a canonical
representation and a classification layer are prerequisites, not refinements. Third,
techniques developed for SQL DBMSs assume a single language and a single query planner and do
not transfer to lazy/eager DataFrame execution or Arrow in-memory layout. Fourth, a fuzzing
finding is not a bug: it becomes a contribution only after it is re-checked, minimized,
deduplicated by root cause, and independently confirmed upstream.

\subsection{Gap}

Prior work has advanced three strands separately. DBMS logic-bug testing provides powerful
oracle families---pivoted query synthesis, non-optimizing reference engines, query
partitioning, cardinality and metamorphic oracles---but targets SQL engines. DataFrame and
workload-generation work targets Python data-processing libraries, but emphasizes workload
realism and replay rather than a cross-model semantic oracle. Specialized data-system testing
targets graph, spatial, or time-series engines. To our knowledge no prior published system
unifies DataFrame APIs, Arrow processing, embedded SQL, and analytical query engines under one
typed testing surface with one normalizer, one compositional oracle, and one evidence
pipeline; we deliberately do not claim priority for differential, metamorphic, or DataFrame
testing as techniques.

\subsection{Thesis}

A single typed workflow IR with capability-aware lowering, a shared semantic normalizer, a
compositional multi-endpoint oracle complex, and an end-to-end evidence pipeline can find,
reproduce, and substantiate latest-version silent wrong-result bugs across the whole tabular
data-processing ecosystem, and can do so with controlled noise and auditable provenance.

\subsection{Contributions}

\begin{itemize}[leftmargin=*,itemsep=2pt]
  \item \textbf{C1 (unified testing surface).} A typed workflow IR and capability-aware
  lowering that place Python DataFrame APIs, Arrow compute, embedded SQL, and query engines on
  one testing surface (\S\ref{sec:approach-surface}), evaluated for reuse and per-family
  coverage in RQ6.
  \item \textbf{C2 (normalizer + compositional oracle).} A semantic normalizer and a
  multi-endpoint oracle complex that fuses differential checking, metamorphic relations, and
  deterministic latest-version probes, with derivation, applicability, and observation
  certificates and separate \textsc{satisfied}/\textsc{violated}/\textsc{inapplicable}/
  \textsc{inconclusive} verdicts (\S\ref{sec:approach-oracle}); evaluated for noise control
  (RQ2) and oracle complementarity (RQ3).
  \item \textbf{C6 (evidence pipeline).} An end-to-end pipeline from candidate to
  issue-ready root to independently confirmed family, with fresh/replay physical separation
  and a strict counting contract (\S\ref{sec:approach-evidence}); evaluated in RQ5.
  \item \textbf{Oracle core (C7).} A finite compositional semantic HyperContract compiled by
  forward-property inference and backward-observability analysis, driving a coverage-debt
  scheduler over a declared target lattice, with authority-bound replayable evidence
  (\S\ref{sec:approach-hypercontract}, \S\ref{sec:impl-authority}); evaluated in RQ-C7.
\end{itemize}

\subsection{Results preview}

Using \sys{} across DataFusion, Polars, DuckDB, and PyArrow, we report \textbf{nine strict
confirmed latest-version bug families} (seven upstream-labeled, two fixed upstream), with one
further finding pending independent confirmation and one matching a pre-existing upstream
issue reported separately (Table~\ref{tab:bugs}). \todo{Add regenerated RQ2/RQ3/RQ5/RQ6
summary sentences once W3 completes.}

\subsection{Artifact availability}

\todo{Describe the artifact package: frozen source commit, exact environment lock, seed
policy, raw journals, receipts, and the single reproduction script.}

\section{Background and Motivating Examples}
\label{sec:background}

\subsection{The heterogeneous tabular stack}

We group targets into four backend families:
\emph{DataFrame APIs} (pandas, Polars eager/lazy/streaming),
\emph{Arrow processing} (PyArrow tables, record batches, compute kernels, layout variants),
\emph{embedded SQL} (DuckDB, SQLite), and
\emph{query engines} (DataFusion, chDB).
Table~\ref{tab:targets} lists the families, execution models, and the capability axes along
which a workflow is lowered. \todo{Generate Table~\ref{tab:targets} from
\texttt{datadiff targets} / \texttt{semantic-registry}.}

\subsection{Silent wrong results}

\todo{Write 2--3 concrete worked examples drawn from the nine strict confirmed families, each
with: the workflow, the two backends, the expected vs observed result, the root cause, and the
upstream issue URL. Use \texttt{experiments/latest_confirmations.json} as the source of truth.}

\begin{table}[t]
\caption{Strict confirmed latest-version bug families by backend.
Source: \texttt{experiments/latest\_confirmations.json}; see
\texttt{paper/CLAIMS\_RECONCILIATION.md}.}
\label{tab:bugs}
\small
\begin{tabular}{@{}llr@{}}
\toprule
Backend family & Root cause & Count \\
\midrule
DataFusion & grouped top-$k$ null sort key; group-by after nested order/limit; negative zero;
distinct null top-$k$; limit idempotence & 5 \\
Polars & reflected arithmetic operand order; group-by aggregation (fixed) & 2 \\
PyArrow & sliced boolean group-by \texttt{any}/\texttt{all} & 1 \\
DuckDB & semi/anti-join rewrite (fixed) & 1 \\
\midrule
\textbf{Strict confirmed} & & \textbf{9} \\
\bottomrule
\end{tabular}
\end{table}

\subsection{Oracle families in testing}

Differential testing compares multiple implementations on the same input; metamorphic testing
checks relations preserved by semantics-preserving transformations; deterministic probes encode
known high-value invariants. In DBMS testing these appear as pivoted query synthesis (PQS),
non-optimizing reference engine construction (NoREC), query partitioning (TLP), equivalent
expression transformation, and constraint-based oracle construction. \todo{Add verified
citations from \texttt{refs.bib}; verify venue/year per entry.}

\subsection{Why a shared semantic layer is required}

A DataFrame API returns an index and a dtype alongside values; Arrow returns a chunked,
possibly sliced column; SQL returns an unordered multiset; a query engine returns an execution
plan with optimizer-dependent ordering. Comparing these outputs without canonicalization
yields false positives on presentation differences and false negatives on real semantic
differences. \S\ref{sec:approach-normalizer} describes the normalizer that makes the
comparison meaningful.

\section{Approach}
\label{sec:approach}

\subsection{Overview}

Figure~\ref{fig:overview} shows the pipeline:
typed workflow generation $\rightarrow$ capability-aware lowering $\rightarrow$ parallel
execution over backend families $\rightarrow$ semantic normalization $\rightarrow$
compositional oracle evaluation $\rightarrow$ triage/classification $\rightarrow$ evidence
pipeline. \todo{Create Figure~\ref{fig:overview} (TikZ or PDF) showing the six stages and the
fresh/replay separation boundary.}

\subsection{Unified testing surface}
\label{sec:approach-surface}

\todo{Describe: (i) the typed workflow IR and small expression AST; (ii) the target registry
that declares each backend's family and capabilities; (iii) capability-aware lowering that
maps one workflow onto pandas/Polars APIs, Arrow compute, SQL, and query-engine execution;
(iv) the shared operation- and case-level semantics that keep generator, scheduler, oracle and
triage from drifting. Point to \texttt{src/datadiff/ccs_ir.py}, \texttt{dsl.py},
\texttt{targets.py}, \texttt{adapters/}, \texttt{backends/}.}

\subsection{Semantic normalization}
\label{sec:approach-normalizer}

\todo{Describe the canonical form: dtype promotion and canonical dtype names, null/NaN/negative
zero equivalence rules, index and column-order erasure where semantics allow, row ordering and
tie determinism, Arrow layout erasure (slice offset, chunking, dictionary encoding), numeric
tolerance policy, and error/timing classification. State explicitly which differences are
\emph{semantic} and which are \emph{presentational}. Point to \texttt{normalizer.py} and
\texttt{canonicalization.py}.}

\subsection{Compositional multi-endpoint oracle}
\label{sec:approach-oracle}

\todo{Describe: differential oracle across backend families; metamorphic relations on
semantics-preserving transformations (sort, filter/take, rebuild, pagination, string
normalization idempotence, coalesce idempotence, semi/anti join with duplicate keys); and
deterministic latest-version probes. Explain per-case fusion and confidence, the four verdicts,
and root-cause classification. \textbf{Recount probes and relations from the live registry
before publishing any number} (see \texttt{paper/CLAIMS\_RECONCILIATION.md} \S3).}

\subsection{Compositional semantic HyperContract (oracle core)}
\label{sec:approach-hypercontract}

\todo{Describe the phase-6 design: finite multi-endpoint contracts; forward semantic-property
inference and backward observability analysis; derivation / applicability / observation
certificates per observed obligation; staged planner vs authority exact comparator with
escalation; the declared target lattice (16 fresh families / 232 cells, 384 single-axis
contrast edges, 502 target-control pair obligations, 9 regression families / 144 cells); and
the coverage-debt scheduler with deterministic seed substreams. This is the paper's
software-engineering depth; keep it factual and separate unrun gates from completed ones.}

\subsection{Search and scheduling}

\todo{Describe typed generation profiles, mutation operators, the multi-scope adaptive
selection across case/profile/target/backend dimensions, the quality-diversity seed archive,
and lane scheduling. Frame as the mechanism behind RQ4; do not claim priority.}

\subsection{Evidence pipeline and counting contract}
\label{sec:approach-evidence}

\todo{Describe the five counting tiers (finding, recheck survivor, native candidate,
issue-ready root, strict confirmed family), family key = root cause + suspicious backends,
fresh/replay physical separation, reduction, issue bundling, and confirmation tracking. Point
to \texttt{triage.py}, \texttt{family\_novelty.py}, \texttt{reducer.py},
\texttt{issue\_bundle.py}, \texttt{final\_readiness.py}.}

\section{Implementation and Reproducibility Discipline}
\label{sec:implementation}

\subsection{Targets, adapters, and capability matrix}

\todo{List implementations and LOC for the target registry, adapters per family, and the
capability matrix; report adapter LOC and time-to-add-a-target as RQ6 evidence.}

\subsection{Determinism and parallel execution}

\todo{Describe deterministic seed substreams, worker-count invariance, retry/completion-order
invariance, and the cache-key boundary that includes evidence tier. Point to
\texttt{execution\_core}, \texttt{parallel}, and the parallel-scaling producers.}

\subsection{Authority-bound evidence}
\label{sec:impl-authority}

The phase-6 work adds an explicit evidence-authority layer: a dynamic plan with frozen seed
blocks, work units, denominators, resource limits, and command digests; an append-only typed
artifact--receipt index in which every receipt binds source snapshot, authority, task/result
identity, seeds, outcomes, raw artifact hashes, and transcript; and fail-closed admission in
which wrong source, substituted identity, duplicate digest, stale authority, malformed raw
artifact, or diagnostic-only reference is rejected as \texttt{not\_evaluated}. This layer is
itself an engineering contribution: it makes each reported number replayable from durable raw
bytes and prevents diagnostic artifacts from earning gate credit. \todo{Describe the ten
producer kinds and the plan/index verifiers; report current completion status honestly, since
the gates are not yet green (\texttt{paper/research/CODE\_GAP\_REPORT.md}).}

\subsection{Reproduction package}

\todo{Specify: pinned source commit, environment lock with hashes, the single reproduction
script, expected wall time, and the exact commands for every table and figure.}

\section{Evaluation}
\label{sec:eval}

\subsection{Experimental setup}

\todo{State machine (CPU/RAM), frozen environment with package hashes, source commit, backend
versions, seed policy (excluded blocks; fresh pair 30700001/30700102), lane specification,
case budget, timeouts, and the counting contract. Reference
\texttt{paper/EXPERIMENT\_MATRIX.md}.}

\subsection{RQ1: Can \sys{} find confirmed latest-version bug families?}

\todo{Fill Table~\ref{tab:rq1} from regenerated campaign manifests. Report strict survivors per
10k cases, unique candidate families, time-to-first, and per-backend strict confirmed families.
State the frozen 9-family baseline (Table~\ref{tab:bugs}) and any new families separately.}

\begin{table}[t]
\caption{RQ1 — fresh discovery per seed block. \todo{regenerate}.}
\label{tab:rq1}
\small
\begin{tabular}{@{}lrrrr@{}}
\toprule
Seed block & Executed & Findings & Survivors & New families \\
\midrule
30700001/30700102 & \todo{} & \todo{} & \todo{} & \todo{} \\
\bottomrule
\end{tabular}
\end{table}

\subsection{RQ2: Does normalization control noise?}

\todo{Report raw findings vs post-classification survivors, FP rate, expected-divergence rate,
and normalization before/after counts for \texttt{legacy\_cartesian} vs
\texttt{contract\_cartesian} vs \texttt{contract\_lattice\_static}.}

\subsection{RQ3: Are the oracle families complementary?}

\todo{Report the union/intersection of families per oracle source, pairwise Jaccard, and the
recounted probe/relation counts.}

\subsection{RQ4: Does typed generation and guidance improve yield?}

\todo{Report valid-program ratio, executed cases/s, candidate families/hour, time-to-first, and
novelty-weighted yield across the adaptive and control arms, with paired bootstrap CIs.}

\subsection{RQ5: Is the evidence pipeline actionable?}

\todo{Report reduction ratio, reproducer success, bundle completeness, readiness, and
flaky/timeout rates over all RQ1 candidates.}

\subsection{RQ6: Is the testing surface transferable?}

\todo{Report workflow reuse rate, capability coverage per family, adapter LOC and
time-to-add-a-target, and per-family confirmed families.}

\subsection{RQ-C7: Does the HyperContract oracle satisfy its gates?}

\todo{Report contract, granularity/reachability, and runtime-correctness gates separately.
Until W2 completes, state clearly that these gates are not yet green and report the current
producer-kind completion count. Do not report partial gate credit.}

\subsection{Baselines}

\todo{Report SQLancer-family and random/unguided comparisons under equal budget and the strict
counting contract; report incompatible baselines separately and never as a win.}

\subsection{Threats to validity}

\todo{Construct, internal, external, conclusion validity; cross-backend divergence is not
automatically a bug; version drift; seed sensitivity; engine coverage imbalance; lost
historical manifests.}

\section{Discussion}
\label{sec:discussion}

\subsection{What transfers across execution models}

\todo{Discuss which parts of the IR, normalizer, and oracle were reusable unchanged across
DataFrame, Arrow, SQL, and query-engine families, and which required family-specific work.
This is the empirical core of the C1 claim.}

\subsection{Engineering lessons}

\todo{Discuss the authority/receipt discipline: what broke when historical manifests were lost,
why physical fresh/replay separation mattered, and what a minimal reproducible evidence layer
requires. Frame honestly; the phase-6 gates are heavy and partly unproven.}

\subsection{Limitations}

\todo{DataFusion-heavy family distribution; confirmation latency and dependence on upstream
maintainers; deterministic probes are manual-knowledge-bound; the HyperContract gates are not
yet green; some claimed mechanisms lack isolating ablations.}

\subsection{Ethics and responsible disclosure}

All reported bugs were submitted upstream through the projects' public trackers; no exploit,
denial-of-service, or data-exfiltration payloads are involved.

\section{Related Work}
\label{sec:related}

\todo{Write four subsections and a comparison table. Verify every citation in \texttt{refs.bib}
before submission.}

\subsection{Testing database engines for logic bugs}
Pivoted query synthesis \cite{rigger2020pqs}, non-optimizing reference engines
\cite{rigger2020norec}, query partitioning \cite{rigger2020tlp}, and the SQLancer tool family
\cite{rigger2020sqlancer}. These provide the oracle vocabulary we adopt, but target SQL DBMSs.

\subsection{Differential and metamorphic testing of data systems}
NEZHA \cite{rui2017nezha} for domain-independent differential testing; SparkFuzz
\cite{sparkfuzz2020} for query-engine correctness regressions; BigFuzz \cite{bigfuzz2020} for
data-analytics framework abstraction; Csmith \cite{yang2011csmith} as the classical
differential-fuzzing reference. \todo{Add metamorphic DB testing (e.g., QTRAN) and recent
equivalent-expression / constraint-based oracles with verified metadata.}

\subsection{DataFrame and workload testing}
FuzzyData and DataFrame workload generation/replay. \todo{Verify metadata; position as workload
generation + replay vs our cross-model semantic oracle.}

\subsection{Search, scheduling, and diversity}
Multi-scope adaptive selection, mutation scheduling (MOPT \cite{mopt2019}), and quality-diversity
archives (MAP-Elites \cite{mouret2015mapelites}). \todo{Position as borrowed mechanisms with
new scope unification, not as primary novelty.}

\begin{table}[t]
\caption{Positioning. \todo{Complete from verified related work:
system, target models, oracle families, evidence pipeline, cross-model reuse.}}
\label{tab:related}
\small
\begin{tabular}{@{}lcccc@{}}
\toprule
System & Models & Oracles & Evidence pipeline & Cross-model reuse \\
\midrule
PQS/NoREC/TLP & SQL & 1 & partial & no \\
FuzzyData & DataFrame & \todo{} & \todo{} & no \\
SparkFuzz & query engine & differential & no & no \\
\sys{} & DataFrame+Arrow+SQL+QE & multi-endpoint & yes & yes \\
\bottomrule
\end{tabular}
\end{table}

\section{Conclusion}
\label{sec:conclusion}

\todo{Two paragraphs. Restate the problem and thesis; summarize C1/C2/C6 and the oracle core;
state the nine strict confirmed families and the regenerated RQ findings; close with the
methodological lesson that cross-ecosystem tabular correctness needs a shared semantic testing
surface and an auditable evidence pipeline, not another single-engine oracle.}

\section*{Data Availability}
\todo{Point to the artifact package: frozen commit, environment lock, seed policy, raw
journals, receipts, and the single reproduction script.}


\bibliographystyle{ACM-Reference-Format}
\bibliography{refs}

\appendix
\section{Appendix: Artifact and Additional Results}
\label{sec:appendix}

\subsection{A. Reproduction}

\todo{Exact commands, expected wall time, and per-table mapping.}

\subsection{B. Confirmed bug families}

\todo{Full table with backend, root cause, upstream issue URL, state, labels, and the minimal
reproducer path for all 9 strict + 1 pending + 1 non-original.}

\subsection{C. Seed policy and exclusions}

\todo{Excluded smoke/yield-gate blocks; longhaul reservation; fresh pairs used.}

\subsection{D. Additional ablations}

\todo{Supplementary arms, negative results, and zero-yield runs.}


\end{document}
